\chapter{Project Requirements Specification}
\label{ch:prs}

\section{Purpose}

This chapter specifies what the OPG analysis system must do and the conditions under which it must
do it. It is written for the project guide, the examiners and any developer who maintains the code.
The requirements describe the system as delivered in Phase 3; Chapter~\ref{ch:results} reports how
well each measurable requirement is met.

\section{Product Scope}

The product is a decision-support prototype that takes one panoramic radiograph and returns a
structured, per-tooth report: which teeth are present, their FDI numbers, whether each shows caries
and how deep, and whether each overlaps an expert-style abnormality. It is meant for demonstration,
teaching and research. It is not a medical device, makes no diagnosis, and states so in every report.

\section{Product Perspective}

The system is self-contained (Figure~\ref{fig:context}). A user supplies an image through a command
line or a browser; the system loads hash-pinned model checkpoints from local storage, runs on the
CPU, and writes a JSON report and an optional overlay. It needs no network connection at inference
time, stores no patient data, and has no external dependencies beyond the Python packages it ships
with.

\begin{figure}[!htbp]
\centering
% System context. Components traced to: scripts/run_inference.py (CLI, load_models, run_pipeline),
% app.py (Streamlit), Dockerfile, models/README.md (checkpoint store), paper/paper.md Table 3 (datasets).
\begin{tikzpicture}[node distance=7mm and 10mm]
\node[stage,text width=50mm,minimum height=13mm] (sys) {OPG analysis system\\[-1pt]{\normalfont six-stage pipeline, CPU}};
\node[soft=p2grey,text width=50mm,below=4mm of sys] (store) {Model store: seven hash-pinned checkpoints (313 MB)};
\begin{scope}[on background layer]
\node[draw=p2indigo,dashed,rounded corners=4pt,line width=0.8pt,fit=(sys)(store),inner sep=6pt,
      label={[lbl,anchor=south]north:system boundary}] (bd) {};
\end{scope}
\node[io,text width=30mm,left=22mm of sys,yshift=10mm] (den) {Dentist / student\\[-1pt]{\normalfont browser}};
\node[io,text width=30mm,left=22mm of sys,yshift=-10mm] (res) {Researcher\\[-1pt]{\normalfont command line}};
\node[rept,text width=30mm,right=22mm of sys,yshift=10mm] (rep) {JSON report\\[-1pt]{\normalfont \texttt{opg-report/1.6}}};
\node[rept,text width=30mm,right=22mm of sys,yshift=-10mm] (ovl) {Overlay image};
\node[soft=p2teal,text width=44mm,below=10mm of store] (data) {Public datasets: DENTEX, Tufts, CariesXrays, Kaggle (training only, offline)};
\node[soft=p2maroon,text width=40mm,left=8mm of data] (ext) {Locked external cohorts: Ningxia, UFBA-425 (evaluation only)};
\node[soft=p2orange,text width=36mm,right=8mm of data] (dev) {Developer: tests, Docker build};
\draw[arr] (den.east) -- node[lbl,above,sloped]{upload OPG} (sys.west|-den.east);
\draw[arr] (res.east) -- node[lbl,below,sloped]{\texttt{-{}-image}} (sys.west|-res.east);
\draw[arr] (sys.east|-rep.west) -- (rep.west);
\draw[arr] (sys.east|-ovl.west) -- (ovl.west);
\draw[darr] (data.north) -- node[lbl,right]{trained checkpoints} (store.south);
\draw[darr] (ext.north) |- ($(store.south west)+(0,-2mm)$);
\draw[darr] (dev.north) |- ($(store.south east)+(0,-2mm)$);
\end{tikzpicture}

\caption[System context diagram]{System context. Users and data sources outside the boundary; the
pipeline, its model store and its outputs inside.}
\label{fig:context}
\end{figure}

\section{User Classes and Characteristics}

The system serves four classes of user, each through a different interface
(Table~\ref{tab:users}).

\begin{table}[!htbp]
\centering
\caption[User classes]{User classes and how each uses the system.}
\label{tab:users}
\small
\begin{tabularx}{\textwidth}{@{}L{3.4cm}YL{4.2cm}@{}}
\toprule
User class & Needs & Interface\\
\midrule
Dentist or dental student & Upload an OPG, see the numbered overlay and the per-tooth findings, read the limitations & Streamlit web demo\\
Researcher & Run the pipeline on many images, change checkpoints, reproduce a reported number & Command-line tool, evaluation scripts\\
Developer / maintainer & Retrain a stage, add a test, rebuild the container & Source code, test suite, Dockerfile\\
Examiner & Check that a number in the report traces to an artefact & Number-verification script, saved outputs\\
\bottomrule
\end{tabularx}
\end{table}

\section{Operating Environment}

Inference runs on a laptop or desktop CPU under Windows, macOS or Linux with Python 3.12, or inside
the provided Docker image (a CPU-only build on \texttt{python:3.12-slim}). Training was done on
cloud and laptop accelerators and is not needed to use the system.

\section{Functional Requirements}

The functional requirements, and the file that implements each, are listed in
Table~\ref{tab:fr}.

\begin{table}[!htbp]
\centering
\caption[Functional requirements]{Functional requirements and where each is implemented.}
\label{tab:fr}
\small
\begin{tabularx}{\textwidth}{@{}L{1.4cm}YL{4.0cm}@{}}
\toprule
ID & Requirement & Implemented in\\
\midrule
FR-01 & Accept one OPG as PNG or JPEG through a command-line argument or a browser upload. & \texttt{run\_inference.py}, \texttt{app.py}\\
FR-02 & Check the input and record warnings (for example an unusual bit depth) in the report rather than failing silently. & \texttt{check\_input}\\
FR-03 & Detect every tooth and return a box and a confidence per tooth, after correcting the image orientation. & Stage 1, \texttt{detect\_upright}\\
FR-04 & Assign one FDI code to each detected tooth, at most once per code, with a confidence and flags for doubtful codes. & Stage 2, Hungarian assignment\\
FR-05 & Find caries lesions on the whole image and attach each lesion to the tooth that contains most of it. & Stage 3a, lesion detector\\
FR-06 & Offer the retained per-crop caries classifier as an alternative method. & \texttt{-{}-caries-method gate}\\
FR-07 & Grade each caries-positive tooth as caries or deep caries. & Stage 3b\\
FR-08 & Flag each tooth whose expanded crop overlaps an expert-style abnormality; withhold the severity grade unless requested. & Stage 4a, \texttt{-{}-bone-severity-grade}\\
FR-09 & Mark each tooth assessable or not, with reasons, when no code was assigned or the numbering confidence is below 0.30. & Stage 5\\
FR-10 & Write one JSON report with a summary, per-tooth records, model metadata and a limitations block. & Stage 6, schema \texttt{opg-report/1.6}\\
FR-11 & Optionally write an annotated overlay image. & \texttt{-{}-overlay}, \texttt{draw\_overlay}\\
FR-12 & In the web demo, show the overlay, a per-tooth table, the limitations and the model versions, and offer the JSON report for download. & \texttt{app.py}\\
FR-13 & Let the user select any retained checkpoint to reproduce an earlier result. & CLI flags (Appendix~B)\\
FR-14 & Record provenance in every report: checkpoint paths and SHA-256 hashes, code version and device. & \texttt{runtime\_info}\\
\bottomrule
\end{tabularx}
\end{table}

\section{Non-Functional Requirements}

The non-functional requirements are listed in Table~\ref{tab:nfr}.

\begin{table}[!htbp]
\centering
\caption[Non-functional requirements]{Non-functional requirements, with the measure used to check each.}
\label{tab:nfr}
\small
\begin{tabularx}{\textwidth}{@{}L{1.7cm}L{3.0cm}YL{3.3cm}@{}}
\toprule
ID & Quality & Requirement & Measure\\
\midrule
NFR-01 & Performance & Report one OPG in a few seconds on a laptop CPU. & 1.70 s per image measured (Section~\ref{sec:res-runtime})\\
NFR-02 & Hardware independence & Need no GPU at inference. & CPU-only container\\
NFR-03 & Reproducibility & Pin every default checkpoint by hash; fix random seeds; version the report schema. & SHA-256 pins, seed 42\\
NFR-04 & Traceability & Trace every reported number to a saved artefact. & 215 numbers asserted\\
NFR-05 & Reliability & Refuse to fuse misaligned stage outputs; guard against NaN confidences. & Assertions in \texttt{build\_report}\\
NFR-06 & Usability & One command, or one browser upload; overlay colours distinguishable under red--green colour blindness. & Grey, orange and blue overlay\\
NFR-07 & Safety and honesty & State in every report that it is not clinically validated, and show the limitations beside the results. & \texttt{limitations} block\\
NFR-08 & Privacy & Run locally as a single-user demo; store no uploaded image. & No network, no storage\\
NFR-09 & Portability & Run identically from a Docker image. & Image build self-test\\
NFR-10 & Maintainability & Keep stages in separate modules with tests. & 23 test modules\\
NFR-11 & Licence compliance & Show only images whose licence allows it. & DENTEX CC BY-NC-SA 4.0\\
\bottomrule
\end{tabularx}
\end{table}

\section{Data Requirements}

Each stage needs a different label, and no single public dataset carries all of them. The datasets,
their labels and their roles are listed in Table~\ref{tab:datasets}. Two rules apply throughout:
images that are byte-identical to training images are removed from every evaluation set, and the
Ningxia cohort is never used for training or tuning.

\section{Hardware and Software Requirements}

\begin{table}[!htbp]
\centering
\caption[Hardware and software requirements]{Hardware and software requirements for using the system.
Versions: \texttt{requirements-lock.txt} and the Dockerfile.}
\label{tab:hwsw}
\small
\begin{tabularx}{\textwidth}{@{}L{4.0cm}Y@{}}
\toprule
Item & Requirement\\
\midrule
Processor & Any 64-bit laptop or desktop CPU; no GPU needed for inference\\
Operating system & Windows, macOS or Linux; or any host that runs Docker\\
Python & 3.12\\
Deep learning & PyTorch 2.2.0 (CPU build), torchvision 0.17.0, Ultralytics 8.4.127 (YOLO11)\\
Numerical & NumPy 1.26.4 (below 2, which breaks the Hungarian solver import), SciPy 1.12.0, scikit-learn 1.4.0, pandas 2.2.0\\
Imaging & OpenCV, Pillow 10.2.0, Matplotlib 3.11.1\\
Interface & Streamlit 1.64.0 for the web demo\\
Testing & pytest 7.4.4\\
Storage & The default checkpoints, about 0.3 GB, are kept in the repository\\
\bottomrule
\end{tabularx}
\end{table}

\section{Design Constraints}

\begin{itemize}
\item Each stage keeps the crop convention it was trained with: a tight $1.0\times$ crop for caries,
  $1.5\times$ by $1.5\times$ for numbering and $1.5\times$ by $2.0\times$ for the abnormality flag.
  Crossing them produces confident, wrong output with no error (Section~\ref{sec:lld-crops}).
\item The detector input size is fixed at 1,024 pixels, the confidence threshold at 0.25 and the non-maximum suppression (NMS)
  IoU at 0.50, the values the models were trained and evaluated with.
\item Field names in the report are stable across versions, so consumers do not break when a
  checkpoint changes; the meaning of each field is stated in the report itself.
\end{itemize}
% {C091}

\section{Assumptions and Dependencies}

\begin{itemize}
\item Input radiographs are adult panoramic OPGs comparable to the training corpora.
\item The public datasets remain available under their licences.
\item Reported accuracy holds only for cohorts like those evaluated; another population may have a
  different caries prior and therefore a different image-blind floor.
\end{itemize}

\section{Use Cases}

The actors and their use cases are drawn in Figure~\ref{fig:usecase} and described in
Table~\ref{tab:usecases}.

\begin{figure}[!htbp]
\centering
% Use-case diagram. Use cases traced to: app.py (upload, overlay, table, limitations, JSON download,
% model versions), scripts/run_inference.py (CLI flags), scripts/verify_paper_numbers.py, Dockerfile.
\begin{tikzpicture}[
  actor/.style={font=\sffamily\scriptsize,align=center},
  uc/.style={ellipse,draw=p2indigo,line width=0.7pt,fill=p2indigo!6,font=\sffamily\scriptsize,
             align=center,minimum width=40mm,minimum height=8.5mm,inner sep=1pt}]
\newcommand{\stick}[2]{%
  \begin{scope}[shift={#1}]
  \draw[p2grey,line width=0.8pt] (0,0.42) circle (0.13);
  \draw[p2grey,line width=0.8pt] (0,0.29) -- (0,-0.05) (-0.22,0.18) -- (0.22,0.18)
        (0,-0.05) -- (-0.17,-0.35) (0,-0.05) -- (0.17,-0.35);
  \node[actor] at (0,-0.6) {#2};
  \end{scope}}
\node[uc] (u1) at (0,3.0) {Analyse one OPG in the browser};
\node[uc] (u2) at (0,1.9) {View findings and limitations};
\node[uc] (u3) at (0,0.8) {Download the JSON report};
\node[uc] (u4) at (0,-0.3) {Run batch analysis (CLI)};
\node[uc] (u5) at (0,-1.4) {Select a retained checkpoint};
\node[uc] (u6) at (0,-2.5) {Reproduce a reported number};
\node[uc] (u7) at (5.4,-0.9) {Retrain a stage};
\node[uc] (u8) at (5.4,-2.0) {Build and run the container};
\begin{scope}[on background layer]
\node[draw=p2grey,rounded corners=4pt,line width=0.8pt,fit=(u1)(u6)(u7)(u8),inner sep=8pt,
      label={[lbl,anchor=south]north:OPG analysis system}] {};
\end{scope}
\stick{(-4.6,2.2)}{Dentist /\\student}
\stick{(-4.6,-1.0)}{Researcher}
\stick{(-4.6,-3.4)}{Examiner}
\stick{(9.0,-1.5)}{Developer}
\foreach \u in {u1,u2,u3} \draw[p2arrow,line width=0.7pt] (-4.3,2.3) -- (\u.west);
\foreach \u in {u4,u5,u6} \draw[p2arrow,line width=0.7pt] (-4.3,-0.9) -- (\u.west);
\draw[p2arrow,line width=0.7pt] (-4.3,-3.3) -- (u6.west);
\foreach \u in {u7,u8} \draw[p2arrow,line width=0.7pt] (8.7,-1.4) -- (\u.east);
\draw[darr] (u1.south) -- node[lbl,right]{\guillemotleft include\guillemotright} (u2.north);
\end{tikzpicture}

\caption[Use-case diagram]{Use-case diagram of the system.}
\label{fig:usecase}
\end{figure}

\begin{table}[!htbp]
\centering
\caption[Main use cases]{Main use cases.}
\label{tab:usecases}
\small
\begin{tabularx}{\textwidth}{@{}L{3.3cm}L{2.6cm}Y@{}}
\toprule
Use case & Actor & Flow\\
\midrule
Analyse one OPG & Dentist, student & Upload the image in the browser; the system runs all stages and shows the overlay, the per-tooth table and the limitations; the user downloads the JSON report.\\
Batch analysis & Researcher & Run the command-line tool per image (or from a script) and collect the JSON reports.\\
Reproduce a result & Researcher, examiner & Select the checkpoint named in the result's log and run the matching evaluation script; compare with the saved artefact.\\
Retrain a stage & Developer & Rebuild the stage's dataset, train with the recorded recipe and seed, register the new checkpoint's hash, rerun the tests.\\
Deploy the demo & Developer & Build the Docker image; the build fails if a dependency would break the pipeline; run it and open port 8501.\\
\bottomrule
\end{tabularx}
\end{table}
